\documentclass[a4paper,12pt]{report}
\newcommand{\path}{../}
\usepackage[fiche,niveau={1MA1},nfiche={19}, annee={Emilie-Gourd, 2025--2026},
auteur={ns},theme={Fonctions -- Géométrie analytique (suite)}]{\path packages/bandeaux}
\usepackage{\path packages/boites}
\usepackage{\path packages/fontes}
\usepackage{\path packages/courslatex}
\usepackage{\path packages/mathenv}
\allowdisplaybreaks

\begin{document}

\begin{boiteExT}[Déterminer si un quadrilatère est un carré]
    \vspace{18cm}
\end{boiteExT}

\begin{exo}
	Montrer que le quadrilatère $ABCD$ avec $A(0;2), B(1;0), C(3;2)$ et $D(2;3)$ est un carré.
\end{exo}
\begin{exo}
Dans un repère orthonormé, on considère les points $A (0; 4)$, $B (−4; −4)$ et $C (4; 0)$.

\begin{tasks}
\task	Quelle est la nature du triangle $ABC$ ?
\task	Calculer les coordonnées du point D, milieu de $[AC]$.
\task	Prouver que $(BD)$ est la bissectrice de l’angle $\hat{ABC}$.
\end{tasks}

\end{exo}

\begin{exo}
On donne les points $A (−2; −1)$, $B (6; 1)$, $C (5; 5)$ et $D (−3; 3)$ dans un repère orthonormé.

	\begin{tasks}
		\task Prouver que le quadrilatère $ABCD$ est un parallélogramme.
		\task Le parallélogramme $ABCD$ est-il un rectangle ?
	\end{tasks}
\end{exo}
\begin{exo}
	Dans un repère orthonormé, on considère les points $A (−4; 10), B (−2; 2)$, $C (6; 0)$ et $D (4; 8)$.
	\begin{tasks}
		\task Démontrer que $[AC]$ et $[BD]$ ont le même milieu.  Dans la suite de l’exercice, on notera $M$ ce point.
		\task Calculer les valeurs exactes de $AB$ et $BC$.
		\task Déduire des deux questions précédentes la nature du
quadrilatère $ABCD$.
	\end{tasks}
\end{exo}
\end{document}
